% NAACL 2027 Student Research Workshop submission draft (ACL template, review mode).
% Numbers marked \TBD come from the phase-2 evaluation and are filled from eval/results/phase2_summary.json.
\documentclass[11pt]{article}
\usepackage[review]{acl}
\usepackage{times}
\usepackage{latexsym}
\usepackage[T1]{fontenc}
\usepackage[utf8]{inputenc}
\usepackage{microtype}
\usepackage{graphicx}
\usepackage{booktabs}
\usepackage{multirow}
\usepackage{xcolor}
\newcommand{\TBD}[1]{\textcolor{red}{[#1]}}
% Fill in the exact model name and version before submission (see chat; kept out of the repo).
\newcommand{\genmodel}{\TBD{generator model, name and version}}
\newcommand{\judgemodel}{\TBD{second judge model, name and version}}

\title{Content, Not Vocabulary: Evaluating a Style Skill for LLM Prose\\Against Humanizers, Detectors, and Human Essays}

\author{Anonymous submission}

\begin{document}
\maketitle

\begin{abstract}
Advice for making language-model prose read as human usually takes the form of banned-word lists, and commercial ``humanizers'' optimize against AI-text detectors. We test a different approach: a prompt-level style skill that audits content and structure before vocabulary and forbids adding any fact the writer did not supply. Across \TBD{20} writing briefs in five genres (college personal statements and short answers, argument essays, science explainers and literary analysis), we compare five conditions: plain generation, a word-list prompt, a popular open-source humanizer skill, our skill as an editor, and our skill as a writer. Blind panels of model judges ranked our skill's output first in \TBD{x\%} of judgments, and it made \TBD{y} unsupported claims per text against \TBD{z} for plain generation. Three open-source detectors scored every skill-written text below 20\%, but they also flagged real admitted college essays written before ChatGPT: the detector trained on Anthropic-model output scored 0 of 16 such essays below 20\%. Personal statements expose a trade-off between narrative arc and fidelity, which a two-pass drafting rule narrows but does not close. We release the skill, the evaluation data and a coaching mode for application essays.
\end{abstract}

\section{Introduction}

Readers have learned to recognize machine-written prose. Corpus studies trace sharp rises in words such as \emph{delve} and \emph{underscore} after the release of ChatGPT \citep{kobak2025delving,liang2024monitoring}, and instruction-tuned models use present participial clauses at two to five times the human rate \citep{reinhart2025llms}. People who use these tools heavily can identify AI-written articles with near-perfect accuracy, and they cite vocabulary, structure and lack of originality as their evidence \citep{russell2025people}.

Two responses dominate practice. The first is the banned-word list, distributed in prompts, style guides and community ``skills''. The second is the commercial humanizer, which rewrites text until a detector scores it as human \citep{roe2026dramaturgies}. Both treat the surface as the problem. Wikipedia's editors, whose guide most word lists derive from, warn against exactly this: ``Please do not merely treat these signs as the problems to be fixed; that could just make detection harder'' \citep{wikipedia2026signs}. Both responses also ignore a second failure of model prose: models fill gaps with plausible details that nobody supplied \citep{ji2023hallucination}, and in personal writing those details are feelings and memories that belong to someone else.

We built and evaluated an alternative: a style skill (a set of instructions an agent loads when a task matches) that orders its audit from content to staging, inflation, rhythm, tone and only then vocabulary, and that never adds a fact absent from the writer's material. For college application essays it switches to a coaching mode and does not write the essay. This paper reports what happened when we tested it.

Our contributions are:
\begin{itemize}
  \item An evidence-graded design for a writing skill, built from a review of about 40 sources on AI-writing tells and of 15 existing open-source humanizer skills, and released with every intermediate version.
  \item A controlled comparison of five conditions on \TBD{20} briefs, scored by blind judge panels for quality, human-likeness and unsupported claims, and by three open-source detectors.
  \item A detector audit against human writing collected before ChatGPT. On admitted college essays the detectors invert: they flag the human essays and pass the skill's output.
  \item A record of four design iterations on personal statements, including the failures, which shows a trade-off between narrative arc and factual fidelity.
\end{itemize}

\section{Related work}

\paragraph{Stylistic signatures of LLM text.} Excess-vocabulary analyses of PubMed and peer reviews show that LLM adoption shifted the frequency of style words, mostly verbs and adjectives \citep{kobak2025delving,liang2024monitoring,juzek2025delve,geng2024chatgpt}. Beyond words, instruction-tuned models write a noun-heavy, informationally dense register \citep{reinhart2025llms}, reuse syntactic templates from pre-training data \citep{shaib2024syntactic}, and overuse corrective framings such as ``not X, but Y'' even where X is not a live alternative \citep{peng2026notx}. Professional writers editing model paragraphs agree on a taxonomy of idiosyncrasies led by clich\'es and unnecessary exposition \citep{chakrabarty2025salvaged}, and an expert taxonomy of ``slop'' finds that syntax-only features predict slop judgments better than vocabulary \citep{shaib2026slop}. Token-level bans on such patterns degrade output quickly \citep{paech2026antislop}. AI generation leaves a stable stylometric footprint that AI editing of human text does not reproduce \citep{shan2026stylometric}.

\paragraph{Detecting AI text.} Detectors range from trained classifiers to zero-shot statistics \citep{mitchell2023detectgpt,bao2024fastdetectgpt,hans2024binoculars}. The RAID benchmark shows that many are nearly inoperable at low false-positive rates and fail under attacks or unseen generators \citep{dugan2024raid}. Paraphrasing defeats watermark-, classifier- and zero-shot detectors \citep{krishna2023paraphrasing,sadasivan2025reliably}, and simple prompting lowers already modest accuracy \citep{perkins2024simple,weberwulff2023testing}. Detectors misclassify non-native English writing as AI at high rates \citep{liang2023biased} and flag formal human prose \citep{creanga2026paradox}; light, guideline-compliant AI editing is flagged more often than humanizer output \citep{karr2026detection}. Untrained humans detect AI text at about chance \citep{clark2021human,jakesch2023heuristics}, while frequent LLM users do far better \citep{russell2025people}. Text watermarks offer a provenance signal that does not depend on style \citep{kirchenbauer2023watermark,dathathri2024synthid}.

\paragraph{LLMs and admissions essays.} LLM-written admissions essays are linguistically distinct from those of 30,000 applicants and resemble the essays of privileged applicants \citep{lee2025poor,alvero2024hegemony}; essay content and style already tracked household income before LLMs \citep{alvero2021essay}. Estimated LLM use in essays rose sharply in 2024 \citep{lee2026divide}, and in one graduate program, essays flagged as AI-written were associated with lower admission \citep{isley2026admissions}. AI writing help homogenizes writing across authors \citep{padmakumar2024diversity,doshi2024creativity,agarwal2025homogenize} and can shift meaning even when limited to grammar \citep{abdulhai2026distort}, and heavy rewriting erases authorship signals \citep{malik2026erased}.

\paragraph{LLMs as judges.} Strong LLM judges agree with human preferences about as often as humans agree with each other, but show position, verbosity and self-enhancement biases \citep{zheng2023judging,wang2024fair}, and models can recognize and prefer their own outputs \citep{panickssery2024llm}. We discuss how these biases bear on our design in Section~\ref{sec:limitations}.

\section{The skill}
\label{sec:skill}

The released skill is a Markdown instruction file (173 lines) with four reference files: a catalog of tells with fixes, a personal-essay guide, academic-register rules, and before/after examples. An agent loads it when a writing task matches its description. We built it from two reviews: a survey of 15 open-source humanizer skills, which supplied the procedure of \citet{chen2026humanizer} and guardrails from two other MIT-licensed skills, and an evidence-graded catalog of tells assembled from about 40 sources.

\paragraph{Audit order.} The catalog showed that the strongest measured signals are structural (participial tails, corrective framing, importance flags, uniform sentence length) and that expert readers rely on content as much as words \citep{russell2025people,shaib2026slop}. The skill therefore audits a draft in a fixed order: content (the ``portability test'': could this sentence move unchanged into someone else's text?), staging, inflation, stated lessons, rhythm, paragraph endings, tone, vocabulary, and format. A strong tell justifies an edit on one sighting; weak tells, such as a single transition word, only when several cluster.

\paragraph{Hard rules.} The skill never adds a fact, name, number, quote, feeling or memory that the user or the source did not provide; where a sentence needs one, it asks or leaves a bracketed question. It keeps every claim and hedge when editing, uses no evasion tricks (typos, rare synonyms, look-alike characters), and makes minimal edits to human-written text.

\paragraph{Over-correction checks.} Humanizer rewrites create tells of their own: the humanizer's own blind test reports more uniform rhythm and more aphoristic closers after rewriting \citep{chen2026humanizer}. The final check lists the last sentence of every paragraph and rewrites when more than two are short or quotable, and checks that a personal essay still has an arc and a line of reflection.

\paragraph{Modes.} \emph{Write} runs the full procedure; for personal essays it drafts in two passes, first for shape and then as an editor of its own draft (Section~\ref{sec:ps}). \emph{Edit} applies the audit to the user's own draft with minimal changes. \emph{Rebuild}, for drafts the user says an AI wrote, first checks every fact, quote and reading against the user's material; if the draft was built from that material it keeps the checked evidence and writes fresh sentences, and if it was built from the topic alone it treats the draft's examples and plan as the stock answer and writes from the first step. \emph{Detect} quotes passages and names patterns without guessing authorship. \emph{Coach}, used for application essays, interviews the student, proposes a structure, critiques drafts and lists corrections, and does not write the essay or example sentences built from the student's story.

\section{Evaluation}
\label{sec:eval}

We evaluated the skill in three rounds: a development round on 20 briefs, a version comparison on the same briefs, and a held-out round on 16 new briefs written after the final version was fixed. A separate experiment tests the coaching mode.

\paragraph{Briefs.} The development set has 20 briefs in five genres, four each (Table~\ref{tab:briefs}). Four fictional students, each with about 1,100 to 1,600 words of interview notes and answers to the skill's depth questions, supply the material for a Common App personal statement and a 250-word supplemental answer. The other briefs are 11th-grade argument essays and science-club explainers with no facts file, and literary analyses restricted to four passages from a public-domain novel or play, copied from Project Gutenberg and checked by script. The held-out set has 16 briefs: two new students (a personal statement and a short answer each), two arguments, two explainers, two literary analyses, and six genres the skill was never tested on: an email to a teacher, a cold email to a professor, a cover letter, a team blog post, the discussion section of a lab report, and a school news article. Every brief names the only permitted source of specifics; the people in the fact files are fictional.

\paragraph{Conditions.} Fresh \genmodel{} agents wrote every text. \textbf{Plain} (P): the model's default writing. \textbf{Word list} (L): the same request plus an instruction to avoid 30 common AI words and constructions and to vary sentence length. \textbf{Humanizer} (H): the most popular open-source humanizer skill \citep{chen2026humanizer}, applied to P's draft. The skill conditions were the Write mode, the Edit mode applied to P's draft, and from version~5 on a Rebuild mode for drafts the user says an AI wrote. The development round compared P, L, H, Edit and Write (version~4); the version comparison compared P, L, version~4 Write, version~5 Write and version~5 Rebuild; the held-out round compared P, L, H, version~4 Write, version~6 Write and version~6 Rebuild.

\paragraph{Judges.} Judge agents read all texts for a brief in a shuffled order without labels: a strict editor, and the genre's real reader (an admissions officer, an English or science teacher, the recipient of the email, a hiring manager, a team mentor, a newspaper adviser). Each ranked the texts, scored each from 1 to 10 on quality, human-likeness, specificity, voice and structure, rated AI suspicion from 1 to 5, and quoted every claim not supported by the facts file (or every invented statistic, source or factual error when there was none). In the held-out round a third judge, a second genre reader, ran on a different model (\judgemodel{}). We report mean rank with bootstrap confidence intervals, a Friedman test across conditions on per-brief mean ranks, and pairwise Wilcoxon signed-rank tests.

\paragraph{Detectors.} We scored every text with three open-source classifiers: a DeBERTa-v3-large model trained on RAID \citep{desklib2025detector}, a RoBERTa-base model \citep{fakespot2025detector}, and a RoBERTa-large model whose training data includes Anthropic-model output \citep{superannotate2024detector}. Texts are split into chunks of about 350 words, and document scores are length-weighted means. We report the share below 20\%, the threshold below which Turnitin hides its AI score.

\paragraph{Human baselines.} To see how the detectors treat human writing, we scored three sets collected before ChatGPT: 16 admitted college essays published by Hamilton College (classes of 2018 and 2022) and Connecticut College (class of 2025); 18 top-scored argument essays from the PERSUADE 2.0 corpus \citep{crossley2024persuade}; and three science explainers published in 2020.

\paragraph{Coaching experiment.} For each of the six fictional students, an agent playing the student wrote a first draft of the personal statement from their notes, with no style guidance. The draft then received feedback from the skill's Coach mode or from the plain model asked for feedback, and the student agent revised it with each feedback separately. Two blind judges ranked the first draft and both revisions, and a counselor agent rated both pieces of feedback for usefulness and counted the essay-ready sentences each wrote for the student. We also measured how much of each revision was copied from the feedback it received: the share of the revision's word 6-grams that occur in the feedback but in neither the first draft nor the student's notes.

\section{Results}
\label{sec:results}

\subsection{Held-out briefs}

Table~\ref{tab:heldout} gives the main result. On 16 briefs written after the final design was fixed, ten genres among them, the skill's Write mode (version~4) ranked first in 23 of 48 judgments and had a mean rank of 1.92 out of 6, against 5.65 for plain generation (Friedman $\chi^2 = 51.0$, $p < 10^{-8}$). It ranked above plain generation in 47 of 48 judgments, above the word-list prompt in 43 and above the humanizer edit in 44 (Wilcoxon $p \le 0.003$ for each). It made 1.17 unsupported claims per text, against 3.38 for plain generation and 2.98 for the word-list prompt. The Rebuild mode, applied to the plain draft, ranked above the humanizer applied to the same draft in 42 of 48 judgments ($p = 0.0008$) and had the fewest unsupported claims of any condition (0.79).

Version~6, which added rules derived from a diagnosis of the development judgments, did not improve on version~4: it ranked above version~4 in 21 of 48 judgments ($p = 0.35$). Its argument essays, built around a worked numerical case, were called ``arithmetic-heavy for an English class.'' We release version~4's Write procedure with version~6's Rebuild mode added (Section~\ref{sec:skill}).

The third judge, which ran on a different model, gave the same order (version~4 1.94, version~6 2.25, Rebuild 2.44, word list 4.19, humanizer 4.81, plain 5.38), so the ranking does not depend on one judge model, within one model family.

\begin{table}[t]
\centering\small
\setlength{\tabcolsep}{3pt}
\begin{tabular}{lcccc}
\toprule
Condition & Rank & 1st & Human & Unsup. \\
\midrule
Plain & 5.65 & 0 & 4.5 & 3.38 \\
Word list & 4.17 & 3 & 6.5 & 2.98 \\
Humanizer edit & 4.25 & 2 & 5.7 & 1.88 \\
Skill write (v4) & \textbf{1.92} & \textbf{23} & \textbf{7.2} & 1.17 \\
Skill write (v6) & 2.42 & 17 & 7.0 & 0.96 \\
Skill rebuild (v6) & 2.60 & 3 & 7.1 & \textbf{0.79} \\
\bottomrule
\end{tabular}
\caption{Held-out evaluation: 16 briefs, three blind judges each (48 judgments). Rank: mean rank among six texts (1 = best); 1st: first places; Human: judged human-likeness (1--10); Unsup.: claims per text that judges found unsupported by the facts file or wrong.}
\label{tab:heldout}
\end{table}

By genre, version~4 ranked best on personal statements (1.33), arguments (1.33), literary analysis (1.33) and the news article (1.00). It ranked worst among the skill conditions on the blog post (4.67), where the word-list prompt won (1.67), and the word-list prompt also edged it on the cover letter (2.00 against 2.33). With one or two briefs per genre these are descriptive, but both losses are in casual first-person genres, where the skill's audit may strip the looseness the genre wants.

\subsection{Development round and the version comparison}

On the 20 development briefs, version~4's Write mode ranked first in 28 of 40 judgments (mean rank 1.45 of 5; plain 4.62; word list 2.27; humanizer 3.48; Edit mode on the plain draft 3.17). Version~5 added a plain student voice for school essays, a Rebuild mode and a personal-essay checklist. In a new five-way panel on the same briefs it raised judged human-likeness (7.45 against 6.90) but ranked below version~4 in 28 of 40 judgments: judges caught simplified reasoning, and the example stance line in the rule (``I think the policy is wrong'') became the opening of three of four argument essays.

The version comparison also exposed a measurement problem. The plain and version~4 texts were identical in both panels, yet on personal statements version~4 ranked above plain in 8 of 8 judgments in the first panel and in only 3 of 8 in the second. The first panel contained two edits of the plain draft, and judges noticed the shared sentences and marked the whole family down. In the other genres the comparison held (8 of 8 in both panels). Rankings of personal statements depend on what they are ranked against, which is one reason we report the held-out round as the main result.

\subsection{Detectors and human writing}

\TBD{Figure~\ref{fig:detectors}; held-out pass rates per condition.} On the development texts, 17 of 20 texts written with the skill scored below 20\% on all three detectors, against 9 of 20 for plain generation, 11 for the word-list prompt and 10 for the humanizer edit. The advantage came from school genres: every condition passed on personal statements and short answers, including plain generation (scores of 0--3\%), while no plain argument or explainer did. The human baselines went the other way. Of the 16 admitted college essays, 5 scored below 20\% on the first two detectors and none on the third, whose mean for them was 61\% (its scores saturate near 73\%). On personal essays the detectors pass every model-written text we produced and flag most human ones.

\subsection{Personal statements: arc against fidelity}
\label{sec:ps}

Personal statements were the hardest genre throughout development (Table~\ref{tab:iterations}). Plain generation built a better arc by writing connective feelings and causes the student never stated (3 to 11 unsupported claims per essay), and judges preferred it. Forbidding invention flattened the skill's first three versions; version~3, for example, ``walks through the notes in order; climax in paragraph 7 of 9.'' Two-pass drafting in version~4 (write for shape first, then audit the draft as an editor) recovered part of the arc: its first win against plain generation on the original notes, with 4 unsupported claims against 14. Giving both conditions the student's answers to depth questions raised everyone's quality and halved plain generation's invention. On the two held-out students, version~4 ranked first in 5 of 6 personal-statement judgments with 0.67 unsupported claims per essay, against 4.83 for plain generation.

\begin{table}[t]
\centering\small
\setlength{\tabcolsep}{3pt}
\begin{tabular}{llcc}
\toprule
Version & Change & vs.\ plain & Unsup. \\
\midrule
v1 & audit while drafting & 0--4 & -- \\
v2 & stricter sample mode & 0--4 & -- \\
v3 & countable checks & 0--4 & 1 vs 16 \\
v4 & two-pass drafting & 1--3 & 6 vs 19 \\
\bottomrule
\end{tabular}
\caption{Personal-statement iterations in development (pairwise judgments of better writing against plain generation, depth notes for v3--v4; unsupported claims summed over the same pairs). Small samples: 2--4 pairs per cell.}
\label{tab:iterations}
\end{table}

\subsection{Coaching}

Both kinds of feedback improved the first drafts: the draft ranked last in all 12 judgments. The revision made with Coach-mode feedback ranked above the one made with plain feedback in 7 of 12 judgments (not significant), with fewer unsupported claims (1.67 against 2.75) and lower AI suspicion (1.83 against 2.08). The clearer difference was in the feedback itself. The counselor judge preferred Coach mode's feedback for all six students (usefulness 8.67 against 7.17 of 10) and found less invented content in it (1.17 items per student against 5.17). Plain feedback rarely wrote whole sentences for the student (1.33 per response), but it supplied the essay's interpretation: one response handed a student ``the essay's interpretive center'' that he never said, which is the substantive content that admissions policies forbid. The simulated students copied almost nothing from either feedback (under 1\% of word 6-grams), so this experiment cannot measure what a rushed real student would copy. Coach mode's responses were also about 1.5 times longer (2,264 words against 1,468), which may be more than a student will read.

\section{Discussion}

\TBD{Why content-first works: fill from the phase-2 and held-out results.}

\paragraph{Detectors in admissions.} On personal essays the three detectors got the direction wrong: they scored real admitted essays high and model essays written with the skill low. A plausible reason is domain. Two of the detectors were trained on RAID, whose eight domains (scientific abstracts, book summaries, news, poetry, recipes, Reddit posts, reviews and Wikipedia) do not include the personal essay \citep{dugan2024raid}. In most of those domains human writing is uneven and specific and model writing is polished and general. Admitted essays are polished, edited many times, and full of the reflective phrasing that admissions advice recommends, so they look like the model side of the training data; a model told to use concrete, uneven detail looks like the human side. Whatever the cause, a detector score is not evidence that an applicant used AI. Colleges that forbid substantive AI content \citep{commonapp2023fraud,brown2026integrity,yale2026ai} have to enforce it some other way, and our results suggest that style-based detection would fall hardest on applicants who write in the polished register that essay coaching teaches.

\paragraph{Coaching versus ghostwriting.} Personal statements were the one genre where the skill's writing did not clearly win, and the reason is instructive. The essays that judges liked most built an arc, a turning point and a closing insight, and plain generation built them by inventing feelings and causes the notes did not contain. When the skill forbade that, it lost some of the arc. The material that makes an application essay work (what the student actually thought, what they still have not resolved) exists only in the student's head, so the honest way for a model to help is to ask for it. Coach mode does this and leaves the writing to the student, which also keeps the essay inside policies that allow brainstorming and proofreading but not AI-written content \citep{yale2026ai}.

\section{Conclusion}

\TBD{}

\section*{Limitations}
\label{sec:limitations}

\paragraph{Model judges.} All judges are language models from the same family as the generator, and models can recognize and favor their own output \citep{panickssery2024llm}. Self-preference does not by itself favor one of our conditions, since every condition was produced by the same model, but judges may share the generator's taste for a particular kind of polish. We shuffled labels, used two reader lenses, and in the held-out study added a judge running on a different model tier; we did not test judges from other model families or recruit enough human raters. A human rating study with the released survey is the most important missing piece.

\paragraph{Test material.} The student personas are fictional and were written by a model; real interview notes are shorter, messier and less quotable. Coach mode was checked on one real student draft and one pressure test, which shows it can work but not how often it does.

\paragraph{Detectors.} We used three open-source classifiers. Commercial detectors used by schools (Turnitin, GPTZero) were not tested, so the pass rates here do not predict their scores. The human baselines are small (16 admitted essays, 18 PERSUADE essays) and the admitted essays were selected and edited for publication by the colleges.

\paragraph{Scope and development.} One generator model wrote every text, in English, for genres common in US high schools. Version 5 of the skill was revised after the phase-2 results on the same 20 briefs, so its phase-2 scores are in-sample; the held-out study uses 16 new briefs, but they were written by the same team that built the skill. Per-genre results rest on two to four briefs each and should be read as descriptive.

\section*{Ethical considerations}

\paragraph{Dual use.} The skill lowers detector scores as a side effect of better writing, and a student could use it to pass off model-written work as their own. We took three steps. The skill uses no evasion tricks (typos, homoglyphs, invisible characters), so it adds nothing to the attack literature on detectors \citep{krishna2023paraphrasing,sadasivan2025reliably}. For application essays it coaches and does not write, and it states institutional AI policies plainly. And it tells users that recent model output may carry a provenance watermark \citep{bleeping2026watermark}: style editing does not make model text the user's own work.

\paragraph{Detector developers.} Because we report low detector scores for model text and high scores for human essays, we \TBD{notified the developers of all three detectors on <date>}, at least 30 days before any preprint.

\paragraph{Human writing and raters.} We scored published admitted essays and PERSUADE essays and report only their scores; we do not redistribute their texts. No real applicant's essay appears in the released data. Any human rating study uses volunteers who give informed consent, collects no names or contact details, and \TBD{had approval from <school or teacher>}.

\section*{AI use statement}

AI agents (\genmodel{}, Anthropic) did a large share of this work: they surveyed the literature and existing skills, drafted every version of the skill, generated all evaluation texts, acted as judges, wrote the analysis scripts, and drafted this paper. The development log, prompts, scripts and every intermediate version are released. \TBD{The author describes here, accurately, what they did themselves: e.g.\ which decisions they made, which numbers and citations they checked against the data and sources, and which sections they rewrote.}

\bibliography{references}

\appendix
\section{Briefs}
\label{sec:appendix-briefs}
\begin{table}[h]
\centering\small
\caption{Writing briefs.}
\label{tab:briefs}
\TBD{table}
\end{table}

\end{document}
